\section*{Methodology}

I use difference-in-differences to compare treated and control firms before and after Helene. The main regression is:
\begin{equation*}
  Y_{it} = a_i + \ell_t
  + b \cdot \mathrm{Treated}_i \cdot \mathrm{Post}_t
  + g \cdot X_{it} + e_{it}
\end{equation*}
$Y_{it}$ is first R\&D intensity (xrdq/atq) and then cash holdings (cheq/atq). I run the same regression twice, once for each outcome. $a_i$ are firm fixed effects. They control for everything about a firm that does not change over time, like industry or location. $\ell_t$ are quarter fixed effects. They control for shocks that hit all firms in the same quarter. $X_{it}$ are controls: size, leverage, profitability and Tobin's Q. Standard errors are clustered by firm, because the same firm is observed many times.

$\mathrm{Treated}_i$ is 1 if the firm is in a FEMA county for Helene, and 0 if not. $\mathrm{Post}_t$ is 1 from Q4 2024 on, and 0 before. I drop Q3 2024 because Helene happened at the very end of September, so that quarter is neither fully before nor after. The coefficient $b$ is what I care about. It shows how treated firms changed after Helene compared to control firms.

For H1 I expect $b < 0$ when $Y$ is R\&D, and $b > 0$ when $Y$ is cash. If both happen together, it supports the idea that firms cut R\&D to hold more cash as protection. This only works if both groups moved in a similar way before Helene. I will check this with a plot of average R\&D and cash for both groups, and by replacing Post with separate quarter dummies to see if there was any difference before.
